% authority: science_draft
\documentclass[11pt]{article}
\usepackage[margin=1in]{geometry}
\usepackage{array}
\usepackage{longtable}
\usepackage{amsmath,amssymb}
\usepackage[T1]{fontenc}
\usepackage{lmodern}

\newcommand{\EqSrc}{\mathrm{EqSrc}}
\newcommand{\RetainH}{\mathrm{RetainH}}
\newcommand{\GenH}{\mathrm{GenH}}
\newcommand{\EqMScert}{\mathrm{EqMS}^{\mathrm{cert}}_{\mathrm{src}}}
\newcommand{\CTr}{\mathrm{CTr}}
\newcommand{\CInv}{\mathrm{CInv}}
\newcommand{\CFact}{\mathrm{CFact}}
\newcommand{\NarrowMSCertEq}{\mathrm{NarrowMSCertEq}_{v1}}
\newcommand{\RRE}{RR_E}
\newcommand{\MetricData}{\mathrm{MetricData}(E)}
\newcommand{\geff}{g_{\mathrm{eff}}}

\title{V15 P5-T01 Upstream Dependency Audit for \EqSrc, \RetainH, and \GenH}
\author{Ontology Formalizer task overlay for RT-20260703-004}
\date{2026-07-03}

\begin{document}
\maketitle

\section*{Control Status}

This artifact implements the bounded v15 P5-T01 task:
\emph{audit whether current matter-sector continuation depends on
\EqSrc, \RetainH, or \GenH}. It is a draft/control dependency audit only.

It does not adopt a canonical ontology edit, discharge general \EqSrc, adopt
\RetainH, adopt \GenH, adopt a source law, adopt matter semantics, adopt
detector semantics, adopt a coupling law, derive or adopt matter coupling,
construct stress-energy semantics, construct a stress-energy tensor, import a
matter action, derive Einstein equations, promote a benchmark, or claim a
completed derivation.

\section*{Evidence Basis}

\begin{longtable}{p{0.27\linewidth}p{0.66\linewidth}}
\textbf{Tracked source} & \textbf{Dependency relevance} \\
\hline
P2 source-side object/certificate manifest &
Defines declared source-side objects, certificate alternatives
\(\CTr,\CInv,\CFact\), no-target guards, and certificate-indexed
\(\EqMScert\). \\
P2 narrow source-side matter-semantics equivalence theorem &
States and proves \(\NarrowMSCertEq\) only under explicit source certificate
premises and fail-closed missing-certificate branches. \\
P2 Refuter stress artifact &
Confirms that deleting or malformed certificates blocks unconditional
equivalence while preserving the theorem under explicit valid certificates. \\
Gate Chair scoped evidence-status review &
Accepts only the scoped source-extension evidence-status role for the declared
conditional theorem; it does not create reusable source-law authority. \\
P3 certificate algebra primitives and operation laws &
Supply certificate-record primitives and conditional operation laws inside
declared source scope, not general \EqSrc, \RetainH, or \GenH. \\
P4 matter-coupling DAG and semantic-layer note &
Record that matter-sector continuation must not collapse source-side
certificate scope into detector, stress-energy, matter-action, coupling-law,
or Einstein-equation authority. \\
Local EqSrc relabeling isomorphism datum &
Supplies only a local EqSrc datum for a local source pair; it explicitly does
not discharge general \EqSrc, \RetainH, or \GenH. \\
Distance-to-GR ledger &
Records general \EqSrc as missing a general equivalence theorem and records
\RetainH and \GenH as missing primitives. \\
\end{longtable}

\section*{Dependency Vocabulary}

The audit separates three notions.

\paragraph{Certificate-indexed source-side matter-semantics equivalence.}
\(\EqMScert\) is the P2/P3 scoped equivalence relation generated by explicit
valid source certificates over declared source objects. It is not a general
\EqSrc theorem and is not detector, physical, stress-energy, action,
coupling-law, empirical, benchmark, or completed-derivation authority.

\paragraph{General source equivalence.}
\(\EqSrc\) denotes the broader upstream source-equivalence burden: equivalence
under the relevant source variations or source-family transformations. The
tracked local EqSrc datum does not discharge this general burden.

\paragraph{Retention and generation primitives.}
\(\RetainH\) denotes an upstream retention law under the relevant \(H\)-indexed
structure. \(\GenH\) denotes an upstream generator law for the relevant
\(H\)-indexed source family. The Distance-to-GR ledger records both as missing
primitives.

\section*{Plan-Required Audit Questions}

\begin{longtable}{p{0.20\linewidth}p{0.26\linewidth}p{0.22\linewidth}p{0.22\linewidth}}
\textbf{Dependency} & \textbf{P2 theorem answer} &
\textbf{P2 classification} & \textbf{Broader continuation classification} \\
\hline
\(\EqSrc\) &
No. The P2 theorem uses explicit certificate premises
\(\CTr,\CInv,\CFact\) and no-target guards; it does not require a general
\(\EqSrc\) theorem. &
\texttt{not\_required\_for\_current\_scope} &
\texttt{conditionally\_required} for any route that removes explicit
certificate premises or seeks family-wide source equivalence. \\
\(\RetainH\) &
No. The P2 theorem does not quantify over \(H\)-retention and does not use a
retention law. &
\texttt{not\_required\_for\_current\_scope} &
\texttt{conditionally\_required} for any route that must preserve declared
certificate structure under \(H\)-retention. \\
\(\GenH\) &
No. The P2 theorem does not generate source families by \(H\) and does not use
a generator primitive. &
\texttt{not\_required\_for\_current\_scope} &
\texttt{conditionally\_required} for any route that constructs or ranges over
an \(H\)-generated source family. \\
\end{longtable}

None of the three broad upstream dependencies is classified as
\texttt{already\_supplied\_by\_tracked\_source}. The tracked sources supply
explicit certificate records, scoped certificate operation laws, and a local
EqSrc datum, but not a general \EqSrc theorem, not \RetainH, and not \GenH.

\section*{Scoped Avoidance Lemma}

\paragraph{Lemma.}
Let \(A_{\mathrm{src}}\) and \(B_{\mathrm{src}}\) be declared source-side
objects. Suppose a valid explicit source certificate of one of the P2/P3
allowed kinds \(\CTr\), \(\CInv\), or \(\CFact\) relates the declared objects,
the no-target guard passes, and the malformed/missing-certificate fail-closed
branches are inactive. Then the P2 theorem's positive conclusion
\(A_{\mathrm{src}} \,\EqMScert\, B_{\mathrm{src}}\) does not require general
\(\EqSrc\), \(\RetainH\), or \(\GenH\) as a premise.

\paragraph{Reasoning.}
The P2 theorem discharges only the declared certificate cases. Its proof
branches over explicit transport, invariance, and factorization certificates.
The P2 Refuter stress artifact shows that removing those certificates blocks
unconditional equivalence rather than repairing the theorem through an
implicit source-equivalence law. P3 supplies only certificate-record operations
inside declared source scope. Therefore the theorem can remain scoped so that
the three upstream primitives are not premises of the present positive branch.

\section*{Broadening Dependency Lemma}

\paragraph{Lemma.}
If a future matter-sector route seeks to remove explicit source certificate
premises, extend the theorem to arbitrary source-family variation, preserve
equivalence under \(H\)-retention, or generate source families by an
\(H\)-indexed rule, then the route becomes conditionally dependent on upstream
\(\EqSrc\), \(\RetainH\), or \(\GenH\), respectively.

\paragraph{Reasoning.}
The present theorem proves only certificate-indexed equivalence. A route that
asks for family-wide equivalence needs a general source-equivalence theorem.
A route that asks for invariance of certificate or semantic status under an
\(H\)-retention operation needs a retention law. A route that asks to construct
or enumerate the relevant \(H\)-indexed source family needs a generator law.
The Distance-to-GR ledger records these as open upstream burdens, so P5-T01
must not silently treat them as supplied.

\section*{Exact Missing Primitives for Broader Routes}

\begin{longtable}{p{0.18\linewidth}p{0.72\linewidth}}
\textbf{Dependency} & \textbf{Missing primitive or theorem if the route broadens} \\
\hline
\(\EqSrc\) &
A general source-equivalence theorem over the relevant declared source-object
and certificate families, including variation hypotheses and no-target
guards. \\
\(\RetainH\) &
A retention law preserving declared source certificate structure, negative
branch behavior, and fail-closed status under the relevant \(H\)-indexed
operation. \\
\(\GenH\) &
A generator law constructing the relevant \(H\)-indexed source family without
target import and with declared-object indexing compatible with the
certificate algebra. \\
\end{longtable}

\section*{Consequence Classification}

P5-T01 does not trigger a freeze for the current P2 theorem, because the
theorem can remain scoped to explicit source certificates and fail-closed
missing-certificate branches. It also does not authorize continuation that
pretends the upstream burdens are already discharged.

The correct consequence is a selector route:
one bounded P5-T02 dependency consequence selector should choose exactly one
next packet among continued scoped matter-sector work with a boundary note, a
narrow \(\EqSrc\) theorem packet, a narrow \(\RetainH\) primitive/theorem
packet, a narrow \(\GenH\) primitive/theorem packet, source-extension
classification enforcement, a protected human-gated ontology route, or a
controlled freeze if repeated missing primitives block the selected route.

\section*{Distance-to-GR Status}

This audit changes no Distance-to-GR ledger row. It only clarifies dependency
status:

\begin{itemize}
\item General \(\EqSrc\) remains open as a general equivalence-theorem burden.
\item \(\RetainH\) remains open as a missing primitive.
\item \(\GenH\) remains open as a missing primitive.
\item Matter coupling remains not derived and not adopted.
\item Einstein equations, benchmark promotion, and completed derivation remain
blocked.
\end{itemize}

\section*{Non-Conclusions}

This audit is not:

\begin{itemize}
\item a general \(\EqSrc\) theorem;
\item a \(\RetainH\) law;
\item a \(\GenH\) law;
\item source-law adoption;
\item matter-semantics adoption;
\item detector-semantics adoption;
\item coupling-law adoption;
\item matter-coupling derivation or adoption;
\item stress-energy semantics, stress-energy tensor, or matter action;
\item an Einstein-equation derivation;
\item benchmark promotion; or
\item completed derivation.
\end{itemize}

\section*{APA 7 Source Materials}

The AEther-Flow Research Project. (2026, June 13). \emph{Local EqSrc
source-relabeling isomorphism datum} [Internal science draft].

The AEther-Flow Research Project. (2026, July 2). \emph{Matter-semantics
equivalence theorem Refuter stress} [Internal science draft].

The AEther-Flow Research Project. (2026, July 2). \emph{Narrow source-side
matter-semantics equivalence theorem} [Internal science draft].

The AEther-Flow Research Project. (2026, July 2). \emph{NarrowMSCertEq Gate
Chair scoped evidence-status review} [Internal scientific gate review].

The AEther-Flow Research Project. (2026, July 2). \emph{Recommendations
implementation plan continue task v15} [Internal implementation plan].

The AEther-Flow Research Project. (2026, July 2). \emph{Source certificate
algebra primitives} [Internal science draft].

The AEther-Flow Research Project. (2026, July 2). \emph{Source certificate
operation laws} [Internal science draft].

The AEther-Flow Research Project. (2026, July 3). \emph{Matter-coupling
dependency DAG v1} [Internal control note].

The AEther-Flow Research Project. (2026, July 3). \emph{Semantic-layer
separation control note} [Internal control note].

\end{document}
